\cotefile{51}
\scanpage{13}
pas toujours ainsi, comme on voit
ainsi. On \deleted{\unsure{peut donc}} suppose $\mathrm{car.}\,k=p>0$,
on prend des ss-schémas finis $\Gamma_1,\Gamma_2$ du
schéma abélien $\deleted{A}B$, tels que $\Gamma'_1\deleted{\unsure{\Gamma_2}}\cap\Gamma'_2=(e)$,
$\Gamma_1\cap\Gamma_2\ne(e)$. On en construit facilement,
p. ex. en partant de $B$ courbe elliptique
telle que $B\otimes k$ n'ait pas de pt d'ordre
$p$ (donc $B\otimes k$ a un seul sous-groupe
fini d'ordre $p$) et $B\otimes K$ ayant deux
sous-groupes distincts d'ordre $p$ (\unsure{l'un|ceci|c'est}
\unsure{faisible|possible|faut-il|faisable} \deleted{\unsure{aux bras}} en \deleted{\unsure{égales car.}}
d'inégales caractéristiques \unsure{?|$0,p$}, \deleted{\unsure{pour que} $\unsure{k/K}$} \inserted{en prenant un}
\deleted{\unsure{pourvu} que $B\otimes K$ ait \unsure{des pts} \inserted{\unsure{de l'ordre}} \unsure{d'ordre} $p$,
on peut \unsure{en trouver} \unsure{beaucoup}, \unsure{sur le}
\unsure{deuxième cas}, \inserted{\unsure{base}} \unsure{en prenant}} \unsure{(à|à} une extension finie de $K$).
On pose $A=B/\Gamma_1\times B/\Gamma_2$, et on y envoie
$B$ diagonalement. OK.
% slanted; long arrow to the space before « Noter »
\inserted{En égales $\mathrm{car.}>0$, on prend pour $B$ un \deleted{\unsure{produit de deux} \unsure{schémas}} schéma abélien \unsure{quelc.|quel.} de dim $\geq2$, \deleted{\unsure{étant la fois}} tel que la puissance $p^{\mathrm{ième}}$ dans $\mathrm{Lie}(B)$ soit nulle, et on fait tourner un sous-espace de dim $2$ ds $\mathrm{Lie}(B)$\unsure{:} \deleted{\unsure{pour}} \inserted{\unsure{des|deux}} \unsure{droites} \unsure{faisant}, \unsure{affines|affirmer|définissant}, \unsure{partir|partant} d'une position où \unsure{ils sont distincts}.}

Noter cependant qu'en car. $0$,
\unsure{comme|pour un} \unsure{tt|$\deleted{A}$|tout} schéma en groupes de type
fini sur un corps est \unsure{séparable|appellé} \unsure{dessus|disons|donc},
% archive sheet number: 12
\scanpage{14}
il est immédiat \unsure{à priori|a priori} que \unsure{tout} ss-schéma
$\deleted{\unsure{A}}$ en groupes \inserted{fermé} $B'$ de $A'=A_K$ est
la \unsure{contraction|restriction} d'un \deleted{\unsure{ab}} ss-schéma abélien
$B$ de $A$. Il s'ensuit, si \unsure{maintt|maintenant}
$A$ est un schéma abélien \deleted{\unsure{pa|projectif}} \deleted{\unsure{quel}} sur un
schéma loc. noeth $S$ \uline{de car. $0$}, \deleted{\unsure{que|tel}} tel
\deleted{le schéma} que $\uuline{\mathrm{Hilb}}_{\inserted{A/S}}\deleted{(A/S)}$ existe \inserted{et est loc. de type fini}
(p. ex. \deleted{\unsure{?}} \deleted{\unsure{que}} \inserted{si} $\deleted{\unsure{B}}A$ projectif sur $S$) alors
le schéma $\unsure{\mathscr{M}|\mathcal{M}}$ sur $S$ de page $2$
\deleted{est tel que} qui classifie les \uline{sous-}schémas \uline{abéliens}
\deleted{(resp. \unsure{séparables})} est
\deleted{tel que ses parties \inserted{\unsure{fermées}} de type fini}
\uline{loc. propre sur $S$}. D'ailleurs, en
toute caractéristique, le $S$-préschéma $N_1$
qui classifie les ss-schémas \deleted{\uline{séparables}} \inserted{en groupes \uline{fermés}} de
$A$ est \deleted{\unsure{localement}} \unsure{sur|\deleted{sur}} \inserted{\uline{loc. non ramifié}} \inserted{\unsure{!|1}} $S$. Donc
En caractéristique $0$ en combinant les,
deux \deleted{\unsure{énoncés}} propriétés, on voit ceci:
% archive sheet number: 13
\scanpage{15}
% slanted
\margin{\uline{N.B.} Il suffit que $A$ soit propre plat \unsure{séparable} sur $S$. (on dit que $A$ est \deleted{\illegible{1}}\inserted{multi}abélien) $\unsure{/}S$ unibranche \uline{suffit}}
\deleted{\uline{Proposition}} \inserted{\uline{Th.}} Soit $A$ un schéma abélien sur le
\deleted{\unsure{$S$ un}} préschéma connexe et \uline{normal} $\unsure{/^{\circ}|/|?}S$ de
car. $0$, de corps $K$. \deleted{\unsure{Pour} \unsure{un}}
\deleted{schéma en groupes \unsure{$L$}} \inserted{Supposons} que
$\uuline{\mathrm{Hilb}}_{A/S}$ existe et soit loc. de type fini sur $S$
[p. ex. que $A$ soit \uline{projectif} sur $S$]. Alors
pour tt sous-schéma en groupes $\deleted{\unsure{G}}B'$ de
$A'=A\otimes K$, il existe un unique sous-schéma
en groupes multiabélien $B$ de $A$,
tel que $B'=B\otimes K$. Si $B'$ est abélien, $B$
est abélien.

\margin{?}Peut-on remplacer l'hypothèse: $S$ normal,
par: $S$ réduit ? \deleted{\unsure{J'en doute fort}} \inserted{Ce n'est pas bien}
probable, mais \deleted{\unsure{sans doute}} c'est OK. \uuline{si} $S$ est géom.
\deleted{\unsure{connexe}}, \inserted{unibranche} \deleted{on pourrait \unsure{quand-même} \unsure{se}
\unsure{prouver|procurer} que les fibres de $B'$\unsure{ont}
\unsure{même} dimension,} \unsure{parce|pour} que $N$ est
\uline{non ramifié} sur $S$.

\uline{Question} \margin{?} Est-il possible \uline{en car. $0$} de définir, si
$A$ est un \deleted{\unsure{sous}} schéma en groupes propre sur $S$,
un plus \unsure{grd|grand} sous-schéma abélien [qui
\unsure{serait|aussi} \unsure{connexe|coïncide} \unsure{qu'on|qd on} \unsure{suppose|regarde} $A$ lui-$=$]. Cas
où $S$ est réduit?
% archive sheet number: 14
\scanpage{16}
\uline{Remarque} Pour $\uuline{\mathrm{Hom}}_{\unsure{S\text{-gr}|S\text{-g}}}(A,B)$ [où $A$ et $B$
sont \deleted{\unsure{abé}}\unsure{]|?}abéliens sur $S$] la situation est analogue.
\uline{sans hypothèse} de \deleted{\unsure{platitude}} caractéristique:
ce schéma sur $S$ [sous réserve d'exister,
et d'être \deleted{de} \inserted{loc. de} type fini] est
loc.? non ramifié et loc. propre sur $S$. D'où
un théorème analogue au précédent sur
l'extension des Hom de schémas multi-abéliens.

[ \uline{Lemme}. Soit $f:X'\unsure{-|\to}X$ un morphisme \uline{fini}
et non \uline{ramifié}, avec $X$ \deleted{normal} loc. noeth.
unibranche. Alors tte section rationnelle de
$X'$ sur $\unsure{x|X}$ est partout définie.

En effet, on est ramené au cas où
\deleted{$f$ est} $X$ est intègre \inserted{affine}, $X'$ intègre, $f:X'\to X$ dominant
fini non ramifié et birationnel, donc
défini par une \unsure{algèbre} $\deleted{\unsure{A}}A'$ sur $A$,
contenue ds la \unsure{cloture|clôture} normale de $A$ dans $K$.
\deleted{\unsure{Comme}} Alors tte fibre a au plus \uline{un}
pt géométrique. Étant séparable, \unsure{l'est|l'on|l'une}
donc \unsure{triv.|iso}, et on finit: $f$ est un isom.].
% archive sheet number: 15
\scanpage{17}
\uline{Proposition} Soient $V$ un anneau de valuation discrète,
de corps $K$, $A_K$ une V.A. définie sur $K$,
$X$ un préschéma projectif sur $V$, et $R_1$ sur $V$,
[\deleted{fibres} plat, fibres simples en codim $\leq 1$],
$u:X_K\to A_K$ une application rationnelle qui
``engendre'' $A_K$. Alors $A_K$ ``se réduit bien''
sur $V$, i.e. il existe un $A$ abélien sur
$V$ tel que $A_K$ soit $K$-isom. à $A\otimes_V K$.

\deleted{On sait que} Soit $C$ une ``courbe
générique'' dans $\unsure{V|X}$, \unsure{universellement|nécessairement|successivement} simple sur
$V$. (On sait que cela existe\ldots). \deleted{Considérons
\unsure{l'hom|l'image}} Il est connu que alors
$C_K\unsure{-|\to}A_K$ est un \uline{morphisme}, et
qui engendre $A_K$ [question géom. sur
un corps alg. clos]. Considérons l'hom
correspondant $\mathrm{Pic}^0(A_K)\to\mathrm{Pic}^0(C_K)$, \deleted{qui
est injectif, \unsure{On|Or} sait que $\mathrm{Pic}^0(\unsure{C_K|A_K})$
``se réduit bien'', \unsure{On}} \inserted{et en \unsure{retransposant|transposant}}
$\mathrm{Pic}^0(C_K)\to A_K$
qui est surjectif [question géom. sur
un corps]. Comme $\mathrm{Pic}^0(C_K)$ se
réduit bien grâce \unsure{à|:} $\mathrm{Pic}^0(C)$, il s'ensuit
\deleted{On} par le th. de Koizumi-Shimura
que \deleted{Pic} $A$ se réduit bien.
% archive sheet number: 16
\scanpage{18}
\uline{Th.} Soient $\deleted{A}$ $V$ anneau de valuation
discrète, $X$ un préschéma projectif sur $V$,
à fibres \deleted{\unsure{simp|simples} séparables} simples en codim $\leq 1$,
\inserted{$X_K$ géom. \unsure{irrédu|irréduc.|irréd.}},
Alors $\mathrm{Pic}^0(X_K)\deleted{\unsure{'|\mathrm{réd}|^{\mathrm{red}}}}$ et $\mathrm{Alb}^0(X_K)$ se réduisent
bien.

Rappelons nous que il existe dans $\mathrm{Pic}^0(X_K)$
un sous-schéma abélien $\mathrm{Pic}^0(X_K)'$, ayant $\unsure{=|\overline{\mathrm{m}}|\text{même}}$
ens. sous-jacent, et que $\mathrm{Alb}^0(X_K)$ est
défini comme le dual de ce dernier.
\deleted{Donc $\mathrm{Alb}^{\unsure{0}}(X_K)$ est isogène à ce dernier,
et la} \margin{$X\unsure{_{K^{\times}1}|_K\times1|_K\,x_1}$}
Faisant une extension non
ramifiée sur $V$, on peut supposer que
$X$ a une section \deleted{\unsure{de}} sur $\mathrm{Spec}(\unsure{W|V})$. Cette
section nous donne un \deleted{\unsure{hom}}
morphisme naturel $X_K\to\mathrm{Alb}^0(X_K)$,
et il est connu que \deleted{\unsure{cet}} ce
morphisme engendre $\mathrm{Alb}^0(X_K)$.
\deleted{\unsure{Le}} (c'est géométrique). Donc par la
proposition précédente, $\mathrm{Alb}^0(X_K)$ se
réduit bien. Mais comme $\mathrm{Pic}^0(X_K)'$
est isogène, il s'ensuit que ce
dernier se réduit bien également.
% archive sheet number: 17
\scanpage{19}
\uline{Théorème}
\margin{\unsure{Grotesque|Grothendieck|Ceci lorsque}, en fait on prouve déjà que $\uline{\mathrm{Pic}}^0_{X/S}$ est \uline{ouvert}, donc en l'occurrence \uline{simple} sur $S$.}
Soient $f:X\unsure{-|\to}S$ un morphisme
% slanted
\margin{Prouvé seulement en \uuline{car. $0$} pour l'instant\ldots.}
projectif, \deleted{plat, \inserted{et} à fibres géométriquement \deleted{\unsure{connexes}} \inserted{intègres} et
simples en codim $1$}, \inserted{\unsure{cf. page|cf. prop.} \unsure{2|?}} la base étant loc.
noethérienne et \uline{réduite}.
% slanted
\margin{est-ce nécessaire ? Non, grâce \unsure{à la simplicité|la rigidité} des Modules formels des V.A. $\ast$}
Alors il existe un
sous-préschéma abélien unique $A$ de $\uline{\mathrm{Pic}}_{X/S}$,
dont l'ensemble sous-jacent soit la réunion
des composantes connexes des fibres de $\uline{\mathrm{Pic}}_{X/S}$.

Soit en effet $P$ la dite réunion. Comme
ses fibres sont \uline{propres} sur $k$ [\uuline{\unsure{réf|vif|vrai}} \ldots],
il s'ensuit \deleted{\unsure{que}} par un théorème connu
[renvoyer aux critères de propreté au Ch. II]
que $P$ est \uline{propre} sur $S$ [question:
en est-il de $=$ de $\mathrm{Pic}^{\tau}_{X/S}$ ? De façon
générale, est-il vrai que $\uline{\mathrm{Pic}}_{X/S}$ est
réunion disjointe d'ouverts \uline{propres} sur
% slanted
\margin{? Non si on suppose $\unsure{X|S}$ seulement formel\ldots.}
$S$ ?]. Nous cherchons un sous-schéma
fermé de $P$, qui soit plat sur $S$ et
qui soit un sous-préschéma en groupes
de $\uline{\mathrm{Pic}}_{X/S}$ \inserted{et abélien sur $S$, \unsure{ensemblt égl à|ensemble ss-jt:|ensemblistement égal à} $\unsure{P|X}$}. \deleted{\unsure{C'est le}} Soit, pour $S'/S$
variable, $F(S')=$ \unsure{ens|les|l'ens.} des sous-schémas
fermés de $P_{S'}$, qui sont des sous-schémas
en groupes de $G$, \inserted{abéliens sur $S$} ensemblistiquement
égaux à $P_{S'}$. Je dis que \unsure{l'on|l'ens} $F(S')$
% archive sheet number: 18
\scanpage{20}
est représentable. En effet, laissant d'abord
tomber la dernière condition, on est
ramené à un sous-schéma \deleted{fermé} de
$\uuline{\mathrm{Hilb}}_{P/S}$ [\inserted{fermé} pour les sous-schémas de $P$
qui sont des sous-groupes de $\mathrm{Pic}$, et
ouvert dans ce dernier pour la condition
abélienne], montrons que la condition
ensembliste est une condition\unsure{,|\,} \uline{ouverte}:
elle est en effet invariante par changement
de base, et si elle est
vérifiée pour un $A\subset P$ en un pt,
elle l'est pour raison de dimensions,
dans un voisinage. De plus,
la projection $M\to S$ de ce problème
modulaire a les propriétés suivantes

(i) les fibres géométriques sont isomorphes \unsure{aux|à un|à}
\deleted{$\unsure{k(x)|k(y)}$} $\mathrm{Spec}\ k(x)$. Évident, en regardant
sur un corps alg. clos, et utilisant \unsure{les|le}
\uuline{th. de rigidité}

(ii) La projection $M\to S$ est propre.

Elle est en effet de type fini\unsure{;|,} séparée, reste
à appliquer le critère valuatif, savoir
ceci: dans le cas où $S$ est le spectre d'un
% archive sheet number: 19
\scanpage{21}
anneau de valuation discrète, et où on se
donne une section \inserted{rationnelle} de $M$ sur $\unsure{B|S}$ i.e. que
on regarde l'unique $A_K$ abélien
sous-jacent à $\mathrm{Pic}^0(X_K)$\deleted{\unsure{,|$\ast$}} \unsure{;|?|$+$} prouver \unsure{qu'en|que}
la section est partout définie. Or par le
\uuline{théorème de Koizumi} $A_K$ ``se réduit bien''
i.e. provient d'un schéma abélien $A$
sur $V$. On va essayer de prolonger
en un \deleted{\unsure{hom}} morphisme $A\to\uline{\mathrm{Pic}}^{-}_{X/S}$
i.e. $A\to P$. C'est possible, grâce au
fait que les fibres de $P$ sur $S$ satisfont
les hypothèses du théorème de \unsure{connexion|connexions} \inserted{\unsure{linéaire|birationnelle}}
de Murre-Chow. [On n'a pas besoin de
ce théorème si $X$ simple sur $S$].
Reste la question si $u:A\to P\deleted{\unsure{x}}$ est une
\uline{immersion} fermée, d'image $P$. L'image a
\deleted{\unsure{est}} évidemment \inserted{\unsure{une fibre|ses fibres}} de dim égale à $\dim A_K$ \deleted{\unsure{génér}}
\inserted{$=\dim\mathrm{Pic}(X_K)$} \deleted{et} [\unsure{continuité} de la dimension des fibres
de $\mathrm{Ker}\,u$] donc de dim égale à
$\dim\mathrm{Pic}(X_k)$ [\uuline{théorème de l'invariance
de la dimension des fibres des Picards}]
% slanted
\margin{Attention, cela est-il vrai si $f:X\to S$ n'est pas simple ??}
% archive sheet number: 20
\scanpage{22}
\mspage{\unsure{11|II}}
Donc l'image de $A$ est $P$. Reste à vérifier
que le noyau de \deleted{$A$ est} $u:A\to P$ \deleted{\unsure{est|$\to$}} [qui, on
le sait, est fini sur $S$, fibre générique étant
le groupe unité] est le groupe unité sur $S$.
En tous cas \deleted{ce groupe est d'ordre} \deleted{\unsure{la fib|Ce fib.}} si
$p$ est l'exposant caractéristique du corps résiduel
$k$, alors la fibre $K_0$ est un $p$-groupe.

\deleted{En effet, si $\xi_0$ est un pt de $K_0$ \unsure{rat./k}
[N.B. on peut supposer $k$ alg. clos \inserted{et $V$ complet}], d'ordre
$n$ premier à $p$, alors $\xi_0$ provient
d'une section $\unsure{\xi|\zeta}$ d'ordre $n$ de $A$ sur $S$,
qui donne une section d'ordre $n$ de
\unsure{le|$\xi$|\&}
\diagram{\begin{tikzcd}[ampersand replacement=\&,column sep=small,row sep=small]
\deleted{Y}\arrow[rr,bend left=13,dash]\arrow[dr] \& \& \deleted{X}\arrow[dl]\\
\&\deleted{S}\&
\end{tikzcd}}
$\uline{\mathrm{Pic}}_{X/S}$ nulle au pt fermé $s$, \unsure{Donc|donc}}

Cela résulte facilement du fait que si
$n$ est premier à $p$, alors $(\uline{\mathrm{Pic}}_{X/S})_n$ est
\deleted{\uline{étale}} \inserted{\uline{non ramifié}} sur $S$, et que \deleted{\unsure{$A$ est|$A_n$}} $A_n$ est \deleted{\unsure{ét|ét.}}
un \uline{revêtement étale} de $S$. Donc si $p=1$
i.e. $k(\unsure{v|V})$ est de car $0$, on a fini.
Sinon, il faut regarder de plus près \unsure{,|.|\, } - - - -

\margin{?}À étudier: noyaux, images d'homomorphismes de
schémas abéliens et multiabéliens, en car $\deleted{\unsure{k|\ne}} \unsure{\ne|k} \unsure{0}$,
notamment sur une base non réduite, ou
% archive sheet number: 21
\scanpage{23}
Schémas abéliens sur un $S$ \uline{régulier}

a) Un schéma préabélien $P$ sur $S$ est projectif sur $S$
(Raisonnement de \unsure{Nakai|Nakaï}-Shimura)

b) Il existe un schéma de translations associé $G$,
et $P$ est un espace \deleted{fibré} principal homogène
\deleted{\unsure{homotopique|homogène}} isotrivial sous $G$.

c) Si $G$ est un schéma abélien sur $S$, alors
$H^1(S,G)\subset H^1(K,G)$. Si $P_K$ est un
fibré principal homogène sous $G_{\unsure{K|k}}$, pour
qu'il \deleted{soit} provienne d'un $P$ pr. hom
ss $G$, il f. et s. qu'il ``se réduise bien''

\margin{\unsure{?}} En tous cas, il y a un plus grand
\deleted{\illegible{1}} ouvert $U$ \deleted{tel} dense de $S$ tel que
$P_K$ provienne d'un $P$ sur $U$ (c'est
l'ens des pts où $P_K$ \deleted{est} se réduit bien)

% slanted
\margin{résulte du th. de pureté}
d) $S-U=Z$ est de codim $1$ en \unsure{tout|tt} point.

Donc
\[
H^1(K/G)/H^1(S,G)\subset
\sum_{\substack{D\ \text{div}\\\deleted{\text{\unsure{irréductibles|irréductible}}}\\\text{premiers}\\\text{de }S}}
H^1(K_{\unsure{D|v}},G)/H^1(\uline{\unsure{O|\mathcal{O}}}_D,\unsure{G}
\]
\margin{\uline{\unsure{gpe|groupe}}: \uline{\unsure{``étudié''|étudié|étudier}} ne \uline{\unsure{ch.ge|change}} \uline{pas par} \uline{complétion}}

Analogue \inserted{th. de pureté} pour la bonne
réduction des schémas abéliens?
% archive sheet number: 22
